\section{Broader Impact}
\label{sec:impact}

Reducing false positives can make security tooling materially more usable, especially for teams that cannot afford dedicated manual triage. At 87.3\% coverage and 89.3\% precision, EVICT would redirect developer attention from roughly 10--20 minutes of manual review per alert to higher-value cases, while returning uncertain alerts to the review queue rather than silently dismissing them. The positive impact is therefore pragmatic and bounded: it reduces the cost of false-positive triage without claiming to replace security expertise.

The primary risk is over-trust. A high-precision triage system can still miss important vulnerabilities if its output is treated as authoritative rather than advisory. Our mitigations are threefold: abstention is a first-class output (not a low-confidence classification), every accepted decision is accompanied by evidence citations and a rationale, and high-severity dismissals are routed to mandatory human audit rather than auto-accepted. These safeguards are specifically designed to prevent the most dangerous failure mode---a critical vulnerability auto-dismissed because the model was confidently wrong---by ensuring that confidence translates to auditable evidence, not just a probability score.

Privacy is a material concern when deployments involve proprietary code. Sending code snippets to public APIs exposes intellectual property and may conflict with enterprise security policies. For such settings, local model deployments or enterprise-hosted APIs are preferable. A third consideration is labor allocation: organizations that treat EVICT's triage assistance as a replacement for security expertise, rather than a tool for focusing that expertise, risk eroding the domain knowledge needed to handle the 10--20\% of alerts that the system correctly defers. Additional discussion of deployment considerations, environmental cost of LLM inference, and dual-use risks is included in Appendix~\ref{app:impact_details}.
